\documentclass[a4paper]{article}

% \usepackage[default]{fontsetup}

\usepackage{fancyhdr}
\usepackage{extramarks}
\usepackage{amsmath}
\usepackage{amsthm}
\usepackage{amsfonts}
\usepackage{tikz}
\usepackage[plain]{algorithm}
\usepackage{algpseudocode}
\usepackage{enumerate}
\usepackage{tikz}
\usepackage{listings}
\usepackage{xcolor}
\usepackage{graphicx}

\usetikzlibrary{automata,positioning}

%
% Basic Document Settings
%  

\topmargin=-0.2in
\evensidemargin=0in
\oddsidemargin=0in
\textwidth=6.5in
\textheight=9.5in
\headsep=0.25in

\linespread{1.1}

\pagestyle{fancy}
\lhead{\hmwkAuthorName}
\chead{\hmwkClass : \hmwkTitle}
\rhead{\firstxmark}
\lfoot{\lastxmark}
\cfoot{\thepage}

\renewcommand\headrulewidth{0.4pt}
\renewcommand\footrulewidth{0.4pt}

\setlength\parindent{0pt}

%
% Create Problem Sections
%

\newcommand{\enterProblemHeader}[1]{
    \nobreak\extramarks{}{Problem \arabic{#1} continued on next page\ldots}\nobreak{}
    \nobreak\extramarks{Problem \arabic{#1} (continued)}{Problem \arabic{#1} continued on next page\ldots}\nobreak{}
}

\newcommand{\exitProblemHeader}[1]{
    \nobreak\extramarks{Problem \arabic{#1} (continued)}{Problem \arabic{#1} continued on next page\ldots}\nobreak{}
    \stepcounter{#1}
    \nobreak\extramarks{Problem \arabic{#1}}{}\nobreak{}
}

\newcommand*\circled[1]{\tikz[baseline=(char.base)]{
		\node[shape=circle,draw,inner sep=2pt] (char) {#1};}}


\setcounter{secnumdepth}{0}
\newcounter{partCounter}
\newcounter{homeworkProblemCounter}
\setcounter{homeworkProblemCounter}{1}
\nobreak\extramarks{Problem \arabic{homeworkProblemCounter}}{}\nobreak{}

%
% Homework Problem Environment
%
% This environment takes an optional argument. When given, it will adjust the
% problem counter. This is useful for when the problems given for your
% assignment aren't sequential. See the last 3 problems of this template for an
% example.
%

\newenvironment{homeworkProblem}[1][-1]{
    \ifnum#1>0
        \setcounter{homeworkProblemCounter}{#1}
    \fi
    \section{Problem \arabic{homeworkProblemCounter} (20\%)}
    \setcounter{partCounter}{1}
    \enterProblemHeader{homeworkProblemCounter}
}{
    \exitProblemHeader{homeworkProblemCounter}
}

\lstset{
    basicstyle=\small\ttfamily,
    numbers=left,               
    numberstyle=\tiny\color{gray},
    keywordstyle=\color{blue},  
    commentstyle=\color{green!50!black}, 
    stringstyle=\color{red},   
    breaklines=true,        
    frame=single            
}


%
% Homework Details
%   - Title
%   - Class
%   - Due date
%   - Name
%   - Student ID

\newcommand{\hmwkTitle}{Homework\ \#01}
\newcommand{\hmwkClass}{Probability \& Statistics for EECS}
\newcommand{\hmwkDueDate}{2024-10-6}
\newcommand{\hmwkAuthorName}{Jianxuan Yao}
\newcommand{\hmwkAuthorID}{2025521024}


%
% Title Page
%

\title{
    \vspace{2in}
    \textmd{\textbf{\hmwkClass:\\  \hmwkTitle}}\\
    \normalsize\vspace{0.1in}\small{Due\ on\ \hmwkDueDate\ at 23:59}\\
	\vspace{4in}
}

\author{
	Name: \textbf{\hmwkAuthorName} \\
	Student ID: \hmwkAuthorID}
\date{}

\renewcommand{\part}[1]{\textbf{\large Part \Alph{partCounter}}\stepcounter{partCounter}\\}

%
% Various Helper Commands
%

% Useful for algorithms
\newcommand{\alg}[1]{\textsc{\bfseries \footnotesize #1}}
% For derivatives
\newcommand{\deriv}[1]{\frac{\mathrm{d}}{\mathrm{d}x} (#1)}
% For partial derivatives
\newcommand{\pderiv}[2]{\frac{\partial}{\partial #1} (#2)}
% Integral dx
\newcommand{\dx}{\mathrm{d}x}
% Alias for the Solution section header
\newcommand{\solution}{\textbf{\large Solution}}
% Probability commands: Expectation, Variance, Covariance, Bias
\newcommand{\E}{\mathrm{E}}
\newcommand{\Var}{\mathrm{Var}}
\newcommand{\Cov}{\mathrm{Cov}}
\newcommand{\Bias}{\mathrm{Bias}}

\begin{document}


% \maketitle
% \thispagestyle{empty}
% \pagebreak

\date{
Due on Sep. 27, 2026, 23:59 UTC+8}
\title{SI 140A.02  Probability \& Statistics for EECS, Fall 2026 \\
Homework 1}
\maketitle
Read all the instructions below carefully before you start working on the assignment, and before you make a submission.
\begin{itemize}
    \item You are required to submit your homework in PDF format generated by \textbf{LaTeX}, otherwise you will get no points of this homework.
    \item You are required to write down all the major steps towards making your conclusions; otherwise you may obtain limited points of the problem.
    \item Write your homework in English; otherwise you will get no points of this homework.
    \item Any form of plagiarism will lead to $0$ point of this homework. 
\end{itemize}
\newpage

\begin{homeworkProblem}[1] 
\textbf{(Story Proof)}
Let \(S(n,k)\) denote the number of ways to partition
\(\{1,2,\ldots,n\}\) into \(k\) nonempty, unlabeled subsets.
Give a story proof for each of the following identities.

\begin{enumerate}
    \item[(a)] Prove that
    \[
        S(n+1,k)=S(n,k-1)+kS(n,k).
    \]

    \item[(b)] Prove that
    \[
        S(n,k)
        =
        \sum_{j=k-1}^{n-1}
        \binom{n-1}{j}S(j,k-1).
    \]
\end{enumerate}

\textbf{Solution:}
\begin{enumerate}
    \item[(a)]
    Pick element $n+1$ out from the set $\{1,2,\ldots,n+1\}$ and consider its placement in the partition.

    Consider 2 cases:
        
    Case 1: The element $n+1$ is in a set by itself. Then the remaining $n$ elements must be partitioned into $k-1$ nonempty, unlabeled subsets. There are $S(n,k-1)$ ways to do this.
    
    Case 2: The element $n+1$ is not a set by itself. Then it must be in one of the $k$ subsets, which are partitioned from the remaining $n$ elements. There are $S(n,k)$ ways to partition the $n$ elements into $k$ nonempty, unlabeled subsets, and for each of these partitions, there are $k$ choices for which subset to place the element $n+1$ into, so there are $kS(n,k)$ ways to do this.

    Add up the two cases, so
    
    $$
    S(n+1,k)=S(n,k-1)+kS(n,k)
    $$

    \item[(b)]
    Pick element $n$ out from the set $\{1,2,\ldots,n\}$ and consider its placement in the partition.
    
    Firstly we devide the remaining $n-1$ elements into two groups: one group of size $j$, and the other group of size $n-1-j$. Group with $n-1-j$ elements are assigned to the same subset, and the group with $j$ elements are partitioned into $k-1$ subsets. There are $\binom{n-1}{j}$ ways to divide the remaining $n-1$ elements into two groups.
    
    Secondly, we need to partition the grouo with $j$ elements into $k-1$ subsets. There are $S(j,k-1)$ ways to do this.

    Then we sum over all possible values of $j$, which can range from $k-1$ to $n-1$, so we have
    
    $$
    S(n,k)
    =
    \sum_{j=k-1}^{n-1}
    \binom{n-1}{j}S(j,k-1)
    $$

    Lastly, we add element $n$ back to the group with $n-1-j$ elements to ensure that all subsets are nonempty. Similarly, $j$ can range from $k-1$ to $n-1$, so we have the final result.

\end{enumerate}

\end{homeworkProblem}
\newpage
\begin{homeworkProblem}[2]

A \emph{nonrepeat password} is a nonempty sequence formed from
\(m\) distinct symbols, where no symbol may appear more than once.
The length of the password is \(\ell\) which could be any integer from \(1\) to \(m\).
Suppose that one password is chosen uniformly at random from all
nonrepeat passwords.

Let \(p_m\) be the probability that the selected password uses all
\(m\) symbols.

\begin{enumerate}
    \item[(a)] Find an expression for \(p_m\).

    \item[(b)] Show that
    \[
        p_m
        =
        \frac{1}{
        \displaystyle\sum_{j=0}^{m-1}\frac{1}{j!}
        }
        \text{ and }
        \lim_{m\to\infty}p_m=\frac{1}{e},
    \]

    where \(e\) is the base of the natural logarithm.
\end{enumerate}

\textbf{Solution:}
\begin{enumerate}
    \item[(a)] 
    The total number of nonrepeat passwords is $\sum_{i=1}^m \binom{m}{i} i!$, and the number of nonrepeat passwords that use all $m$ symbols is $\binom{m}{m} m! = m!$. Therefore,
    
    $$
    p_m = \frac{m!}{\sum_{i=1}^m \binom{m}{i} i!}
    $$
    
    \item[(b)] 
    
    Let $j=m-i+1$, so the expression for $p_m$ can be rewritten as
    
    $$
    p_m = \frac{1}{\sum_{j=0}^{m-1} \frac{1}{j!}}
    $$
    
    As $m \to \infty$, the sum $\sum_{i=0}^{m-1} \frac{1}{i!}$ approaches the Taylor series for $e^1$ ($e^x = \sum_{i=0}^{\infty} \frac{x^i}{i!}$ ), so
    
    $$
    \lim_{m\to\infty}p_m = \frac{1}{e}
    $$

\end{enumerate}

\end{homeworkProblem}

\newpage
\begin{homeworkProblem}[3]
\textbf{(Geometry Probability)}
    You get a stick and break it randomly into three pieces. Two breaking points are chosen independently and uniformly.
What is the probability that you can make a triangle using such three pieces?

\textbf{Solution:}

Let the whole stick be of length 1. Let the two breaking points be $x$ and $y$, where $0 < x, y < 1$.

Assuming $x < y$, the three pieces have lengths $x$, $y-x$, and $1-y$. For these three pieces to form a triangle, the triangle inequality must be satisfied:

$$
\left\{
\begin{aligned}
    x + (y-x) &> 1-y \\
    y + (1-y) &> x \\
    x + (1-y) &> y-x
\end{aligned}
\right.
$$

which simplifies to:

$$
\left\{
\begin{aligned}
    y &> \frac{1}{2} \\
    x &< \frac{1}{2} \\
    x + y &> \frac{1}{2}
\end{aligned}
\right.
$$

It makes a triangle in XOY plane, whose size is $\frac{1}{8}$.

When $x > y$, the probability is the same.

As a result, the total probability is $\frac{1}{4}$.

\end{homeworkProblem}


\newpage
\begin{homeworkProblem}[4]
\textbf{(Birthday Problem)}
In the birthday problem, we assumed that all 365 days of the year are equally likely (and excluded February 29). In reality, some days are slightly more likely as birthdays than others. For example, scientists have long struggled to understand why more babies are born 9 months after a holiday. Let $\mathbf{p}=\left(p_1, p_2, \ldots, p_{365}\right)$ be the vector of birthday probabilities, with $p_j$ the probability of being born on the $j$ th day of the year (February 29 is still excluded, with no offense intended to Leap Dayers). The $k$ th elementary symmetric polynomial in the variables $x_1, \ldots, x_n$ is defined by
$$
e_k\left(x_1, \ldots, x_n\right)=\sum_{1 \leq j_1<j_2<\cdots<j_k \leq n} x_{j_1} \ldots x_{j_k},  \quad 2 \leq k \leq 365.
$$
This just says to add up all of the $\left(\begin{array}{l}n \\ k\end{array}\right)$ terms we can get by choosing and multiplying $k$ of the variables. For example, $e_1\left(x_1, x_2, x_3\right)=x_1+x_2+x_3, e_2\left(x_1, x_2, x_3\right)=x_1 x_2+$ $x_1 x_3+x_2 x_3$, and $e_3\left(x_1, x_2, x_3\right)=x_1 x_2 x_3$ Now let $k \geq 2$ be the number of people.\\
(a) Find a simple expression for the probability that there is at least one birthday match, in terms of $\mathbf{p}$ and an elementary symmetric polynomial.\\
(b) Explain intuitively why it makes sense that $P$ (at least one birthday match) is minimized when $p_j=\frac{1}{365}$ for all $j$, by considering simple and extreme cases.\\
(c) The famous arithmetic mean-geometric mean inequality says that for $x, y \geq 0$
$$
\frac{x+y}{2} \geq \sqrt{x y} .
$$
This inequality follows from adding $4 x y$ to both sides of $x^2-2 x y+y^2=(x-y)^2 \geq 0$ Define $\mathbf{r}=\left(r_1, \ldots, r_{365}\right)$ by $r_1=r_2=\left(p_1+p_2\right) / 2, r_j=p_j$ for $3 \leq j \leq 365$. Using the arithmetic mean-geometric mean bound and the fact, which you should verify, that
$$
e_k\left(x_1, \ldots, x_n\right)=x_1 x_2 e_{k-2}\left(x_3, \ldots, x_n\right)+\left(x_1+x_2\right) e_{k-1}\left(x_3, \ldots, x_n\right)+e_k\left(x_3, \ldots, x_n\right)
$$
show that

\begin{center}
    $P($ at least one birthday match $\mid \mathbf{p}) \geq P($ at least one birthday match $\mid \mathbf{r})$
\end{center}

with strict inequality if $\mathbf{p} \neq \mathbf{r}$, where the given $\mathbf{r}$ notation means that the birthday probabilities are given by $\mathbf{r}$. Using this, show that the value of $\mathbf{p}$ that minimizes the probability of at least one birthday match is given by $p_j=\frac{1}{365}$ for all $j$.

\textbf{Solution:}

\begin{enumerate}
    \item[(a)]
    
    Assume there are $k$ people. Select distinguish $k$ days from 365 days, and assign each person to a day.

    In this case, the probability that there is at least one birthday match is given by

    $$
    k! \cdot p_{j_1} p_{j_2} \cdots p_{j_k}
    $$
    
    and traverse all possible combinations of $k$ days from 365 days, we have

    $$
    P(\text{no birthday match}) = \sum_{1 \le j_1 < j_2 < \cdots < j_k \le 365} k! \, p_{j_1} p_{j_2} \cdots p_{j_k} = k! \, e_k(\mathbf{p})
    $$

    so 

    $$
    P(\text{at least one birthday match}) = 1 - k! \, e_k(\mathbf{p})
    $$

    \item[(b)]
    
    Assume there is an extreme uneven distribution, for example, everyone must have been born on the same day. When $k \ge 2 $, everyone was born on the first day, and the probability of having the same birthday reaches its maximum value of $1 $.

    The most effective way to make these $k$ individuals as scattered as possible is to evenly distribute the probability over 365 days to prevent date collisions. So intuitively, $p_j=\frac{1}{365}$ for all $j$ minimizes the probability of at least one birthday match.

    \item[(c)]
    
    Divide $\{x_1, x_2, \cdots, x_n\}$ into $\{x_1, x_2\}$ and $\{x_3, \cdots, x_n\}$. By the definition, $e_k(x_1, x_2, \cdots, x_n)$ is the sum of all products of $k$ elements from $\{x_1, x_2, \cdots, x_n\}$.
    
    We can divide these products into three categories: those that contain both $x_1$ and $x_2$, those that contain either $x_1$ or $x_2$, and those that contain neither.

    When we select $x_1$ and $x_2$, then we still need to select $k-2$ elements from $\{x_3, \cdots, x_n\}$, which is $e_{k-2}(x_3, \cdots, x_n)$.

    When we select either $x_1$ or $x_2$, then we still need to select $k-1$ elements from $\{x_3, \cdots, x_n\}$, which is $(x_1+x_2)e_{k-1}(x_3, \cdots, x_n)$.

    When we select neither $x_1$ nor $x_2$, then we still need to select $k$ elements from $\{x_3, \cdots, x_n\}$, which is $e_k(x_3, \cdots, x_n)$.

    Add up the three cases, we have

    $$
    e_k(x_1, x_2, \cdots, x_n) = e_{k-2}(x_3, \cdots, x_n) + (x_1+x_2)e_{k-1}(x_3, \cdots, x_n) + e_k(x_3, \cdots, x_n)
    $$

    $r_1+r_2=x_1+x_2$, $r_k = p_k$ for $k \geq 3$, so $e_k(\mathbf{r})-e_k(\mathbf{p})=(r_1 r_2 - p_1 p_2)e_{k-2}(x_3, \cdots, x_n)$

    According to AM-GM inequality, we have $r_1 r_2 \geq p_1 p_2$, so

    $$
    e_k(\mathbf{r}) \geq e_k(\mathbf{p})
    $$

    If and only if $p_1=p_2$, we have $e_k(\mathbf{r}) = e_k(\mathbf{p})$.

    And $P(\text{at least one birthday match} \mid \mathbf{p}) = 1 - k! \, e_k(\mathbf{p})$, so

    $$
    P(\text{at least one birthday match} \mid \mathbf{p}) \geq P(\text{at least one birthday match} \mid \mathbf{r})
    $$

    and when $\mathbf{p} \neq \mathbf{r}$, we have strict inequality.

    $e_k(x_1, \cdots, x_n)$ is a symmetric polynomial, so we can apply the above process to any two elements of $\mathbf{p}$, and we can eventually make all elements equal. Therefore, the value of $\mathbf{p}$ that minimizes the probability of at least one birthday match is given by $p_j=\frac{1}{365}$ for all $j$.

\end{enumerate}

\end{homeworkProblem}

\newpage

\begin{homeworkProblem}[5]

A card collection contains five different types of cards. Each
independent draw produces card type \(i\) with probability \(p_i\),
where
\[
    (p_1,p_2,p_3,p_4,p_5)
    =
    (0.05,0.15,0.20,0.25,0.35).
\]

Let \(Q_n\) denote the probability of collecting all five card types
after \(n\) independent draws. $n \geq 1$.

For each \(i\in\{1,\ldots,5\}\), define
\[
    E_i
    =
    \{\text{card type \(i\) does not appear in the first \(n\) draws}\}.
\]

For any index set
\[
    I\subseteq\{1,\ldots,5\},
\]
define
\[
    E_I=\bigcap_{i\in I}E_i.
\]
Thus, \(E_I\) is the event that none of the card types whose indices
belong to \(I\) appears in the first \(n\) draws. By convention,
\[
    E_{\emptyset}=\Omega,
\]
where \(\Omega\) denotes the entire sample space.

\begin{enumerate}
    \item[(a)]
    Use the inclusion-exclusion principle to prove that
    \[
        Q_n
        =
        \sum_{I\subseteq\{1,\ldots,5\}}
        (-1)^{|I|}
        \left(
        1-\sum_{i\in I}p_i
        \right)^n,
    \]
    where \(|I|\) denotes the number of elements in \(I\).

    \item[(b)]
    Find the smallest integer \(n\) such that $Q_n\geq 0.90.$ You can write some code to calculate it. Include the screenshot of your code and output in the terminal which can support your answer.

    \item[(c)]
    Find the smallest n for which the probability of collecting all five types is at least 0.90 in the uniform case. (No need to include the code)
    Write code to plot $Q_n$ as a function of $n$ for both equal and unequal $p_i$, along with their corresponding probabilities.  (You should incluide the graph.) 
    Compare \(Q_n\) with the corresponding probability between the two cases.
\end{enumerate}

\textbf{Solution:}

\begin{enumerate}
    \item[(a)]
    
    $(1-\sum_{i\in I}p_i)^n$ means the probability of not drawing any card type in $I$ after $n$ independent draws, which can also be interpreted as the probability of drawing at most $5-|I|$ card types in $I^C$ after $n$ independent draws.

    According to the inclusion-exclusion principle, collecting all five card types after $n$ independent draws is equivalent to:

    Include at most 5 card types, exclude at most 4 card types, include at most 3 card types, exclude at most 2 card types, include at most 1 card type, exclude at most 0 card types.

    So when $|I|$ is even, we include, which means we use a coefficient $1$ before it. And when $|I|$ is odd, we exclude, which means we use a coefficient $-1$ before it. Therefore, we have

    $$
    Q_n
    =
    \sum_{I\subseteq\{1,\ldots,5\}}
    (-1)^{|I|}
    \left(
    1-\sum_{i\in I}p_i
    \right)^n
    $$

    \item[(b)]
    
    \begin{lstlisting}[language=Python, caption=Python Code Simulation]
        from itertools import combinations

        p = [0.05, 0.15, 0.20, 0.25, 0.35]

        subsets = []
        for k in range(len(p) + 1):
            for combo in combinations(p, k):
                subsets.append((k, sum(combo)))

        def compute_Q(n):
            return sum(((-1) ** k) * ((1.0 - p_sum) ** n) for k, p_sum in subsets)

        n = 1
        while True:
            q_n = compute_Q(n)
            if q_n >= 0.90:
                print(f"minimal n = {n}, and the corresponding Q_n = {q_n:.6f}")
                break
            n += 1
    \end{lstlisting}

    \begin{lstlisting}[language=Bash, caption=Output]
        minimal n = 46, and the corresponding Q_n = 0.904965
    \end{lstlisting}
    
    \begin{figure}[h]
        \centering
        \includegraphics[width=0.8\textwidth]{Record.png}
        \caption{Output of the Python Code}
    \end{figure}

    \item[(c)]
    
    $$
    \begin{aligned}
    Q_n^{\text{uniform}}
    &= \sum_{k=0}^5 (-1)^k \binom{5}{k} (1 - 0.2k)^n \\
    &= 1 - 5(0.8)^n + 10(0.6)^n - 10(0.4)^n + 5(0.2)^n
    \end{aligned}
    $$

    When $n>19$, $Q_n^{\text{uniform}} > 0.90$.

    \begin{lstlisting}[language=Python, caption=Python Code Plot Generation]
        import matplotlib.pyplot as plt
        from itertools import combinations

        p_unequal = [0.05, 0.15, 0.20, 0.25, 0.35]
        p_equal = [0.20, 0.20, 0.20, 0.20, 0.20]

        def get_subsets(p_list):
            subsets = []
            for k in range(len(p_list) + 1):
                for combo in combinations(p_list, k):
                    subsets.append((k, sum(combo)))
            return subsets

        subsets_unequal = get_subsets(p_unequal)
        subsets_equal = get_subsets(p_equal)

        def compute_Q(n, subsets):
            return sum(((-1) ** k) * ((1.0 - p_sum) ** n) for k, p_sum in subsets)

        n_values = list(range(1, 101))
        Q_unequal = [compute_Q(n, subsets_unequal) for n in n_values]
        Q_equal = [compute_Q(n, subsets_equal) for n in n_values]

        plt.figure(figsize=(10, 6))
        plt.plot(n_values, Q_equal, label='Uniform Case ($p_i = 0.2$)', color='tab:blue', linewidth=2)
        plt.plot(n_values, Q_unequal, label='Unequal Case ($p = [0.05, 0.15, 0.2, 0.25, 0.35]$)', color='tab:red', linewidth=2)

        plt.axhline(0.90, color='gray', linestyle='--', label='Threshold $Q_n = 0.90$')
        plt.scatter([19], [compute_Q(19, subsets_equal)], color='tab:blue', zorder=5)
        plt.annotate(f'Uniform: n=19 (Q≈{compute_Q(19, subsets_equal):.3f})', 
                    xy=(19, compute_Q(19, subsets_equal)), xytext=(22, 0.82),
                    arrowprops=dict(arrowstyle='->', color='tab:blue'))

        plt.scatter([56], [compute_Q(56, subsets_unequal)], color='tab:red', zorder=5)
        plt.annotate(f'Unequal: n=56 (Q≈{compute_Q(56, subsets_unequal):.3f})', 
                    xy=(56, compute_Q(56, subsets_unequal)), xytext=(40, 0.70),
                    arrowprops=dict(arrowstyle='->', color='tab:red'))

        plt.title('$Q_n$ as a Function of $n$ (Uniform vs Unequal Probabilities)', fontsize=14)
        plt.xlabel('Number of Draws ($n$)', fontsize=12)
        plt.ylabel('Probability of Collecting All Types ($Q_n$)', fontsize=12)
        plt.grid(True, linestyle=':', alpha=0.6)
        plt.legend(loc='lower right', fontsize=11)
        plt.tight_layout()
        plt.savefig('contrast_plot.png')
        plt.show()
    \end{lstlisting}

    \begin{figure}[h]
        \centering
        \includegraphics[width=0.8\textwidth]{Code/contrast_plot.png}
        \caption{Contrast Plot of $Q_n$ for Uniform and Unequal Cases}
    \end{figure}

    To any given $n>5$, the probability of collecting all five types in the uniform case is higher than that in the unequal case obviously. What's more, to reach $90\%$ probability, uniform case requires 19 draws, while unequal case requires 56 draws.

    It can be explained by the bottleneck effect, because in the unequal case, the rarest card type (with probability 0.05) is much less likely to be drawn, which significantly slows down the overall collection process;While in the uniform case, all card types have equal probability, making it easier to collect all types in fewer draws.
\end{enumerate}

\end{homeworkProblem}

\end{document}